\renewcommand{\thesection}{\arabic{section}}\renewcommand{\theHsection}{chapter.\arabic{section}}\setcounter{section}{0}\section{极限}
\noindent\begin{minipage}{\textwidth}\ProblemHeading{习题}{1.1}
设 $a$ 是有理数，$b$ 是无理数。求证：$a+b$ 和 $a-b$ 都是无理数；当 $a\ne0$ 时，$ab$ 和 $b/a$ 也都是无理数。
\end{minipage}\par\vfill
\noindent\begin{minipage}{\textwidth}\ProblemHeading{习题}{1.2}
求证：两个不同的有理数之间有无理数。
\end{minipage}\par\vfill
\noindent\begin{minipage}{\textwidth}\ProblemHeading{习题}{1.3}
求证：$\sqrt2$、$\sqrt3$ 以及 $\sqrt2+\sqrt3$ 都是无理数。
\end{minipage}\par\vfill
\noindent\begin{minipage}{\textwidth}\ProblemHeading{习题}{1.4}
把下列循环小数表示为分数：
\SubQ{1}{$0.24999\cdots$。}
\SubQ{2}{$0.\overline{375}$。}
\SubQ{3}{$4.\overline{518}$。}
\end{minipage}\par\vfill
\newpage
\noindent\begin{minipage}{\textwidth}\ProblemHeading{习题}{1.5}
设 $r,s,t$ 都是有理数。求证：
\SubQ{1}{若 $r+s\sqrt2=0$，则 $r=s=0$。}
\SubQ{2}{若 $r+s\sqrt2+t\sqrt3=0$，则 $r=s=t=0$。}
\end{minipage}\par\vfill
\noindent\begin{minipage}{\textwidth}\ProblemHeading{习题}{1.6}
设实数 $a_1,a_2,\ldots,a_n$ 有相同的符号，且都大于 $-1$。证明：\[(1+a_1)(1+a_2)\cdots(1+a_n)\ge1+a_1+a_2+\cdots+a_n.\]
\end{minipage}\par\vfill
\noindent\begin{minipage}{\textwidth}\ProblemHeading{习题}{1.7}
设 $a,b$ 是实数，且 $|a|<1$、$|b|<1$。证明：\[\left|\frac{a+b}{1+ab}\right|<1.\]
\end{minipage}\par\vfill
\noindent\begin{minipage}{\textwidth}\ProblemHeading{例题}{1.8（改写）}
设数列 $a_n=a\ (n=1,2,3,\ldots)$ 为常数列。证明 $\displaystyle\lim_{n\to\infty}a_n=a$。
\end{minipage}\par\vfill
\newpage
\noindent\begin{minipage}{\textwidth}\ProblemHeading{例题}{1.9}
设 $\alpha>0$。证明 $\displaystyle\lim_{n\to\infty}\frac{1}{n^\alpha}=0$。
\end{minipage}\par\vfill
\noindent\begin{minipage}{\textwidth}\ProblemHeading{例题}{1.10}
设数列 $a_1=0.9,\ a_2=0.99,\ a_3=0.999,\ldots$。证明 $\displaystyle\lim_{n\to\infty}a_n=1$。
\end{minipage}\par\vfill
\noindent\begin{minipage}{\textwidth}\ProblemHeading{例题}{1.11}
设 $q$ 是一个给定的实数，满足 $|q|<1$。求证 $\displaystyle\lim_{n\to\infty}q^n=0$。
\end{minipage}\par\vfill
\noindent\begin{minipage}{\textwidth}\ProblemHeading{例题}{1.12}
证明 $\displaystyle\lim_{n\to\infty}\frac{2^n}{n!}=0$。
\end{minipage}\par\vfill
\newpage
\noindent\begin{minipage}{\textwidth}\ProblemHeading{例题}{1.13}
求极限 $\displaystyle\lim_{n\to\infty}\frac{1+2+\cdots+n}{n^2}$。
\end{minipage}\par\vfill
\noindent\begin{minipage}{\textwidth}\ProblemHeading{例题}{1.14}
求极限 $\displaystyle\lim_{n\to\infty}\frac{2^n+5\cdot3^n}{3^n+5\cdot2^n}$。
\end{minipage}\par\vfill
\noindent\begin{minipage}{\textwidth}\ProblemHeading{例题}{1.15}
求 $\displaystyle\lim_{n\to\infty}\sqrt{1+\frac1{n^\alpha}}$，其中 $\alpha$ 是给定的正实数。
\end{minipage}\par\vfill
\noindent\begin{minipage}{\textwidth}\ProblemHeading{例题}{1.16}
求 $\displaystyle\lim_{n\to\infty}\left(\frac1{\sqrt{n^2+1}}+\frac1{\sqrt{n^2+2}}+\cdots+\frac1{\sqrt{n^2+n}}\right)$。
\end{minipage}\par\vfill
\newpage
\noindent\begin{minipage}{\textwidth}\ProblemHeading{例题}{1.17}
设 $a>0$，求证 $\displaystyle\lim_{n\to\infty}\sqrt[n]{a}=1$。
\end{minipage}\par\vfill
\noindent\begin{minipage}{\textwidth}\ProblemHeading{例题}{1.18}
求证 $\displaystyle\lim_{n\to\infty}\sqrt[n]{n}=1$。
\end{minipage}\par\vfill
\noindent\begin{minipage}{\textwidth}\ProblemHeading{例题}{1.19}
求证数列 $\{(-1)^n\}$ 发散。
\end{minipage}\par\vfill
\noindent\begin{minipage}{\textwidth}\ProblemHeading{例题}{1.20（改写）}
证明数列 $1,1,1,2,1,3,\ldots,1,n,\ldots$ 发散。
\end{minipage}\par\vfill
\newpage
\noindent\begin{minipage}{\textwidth}\ProblemHeading{例题}{1.21}
证明如果收敛数列的极限不为 $0$，那么这个数列中至多只有有限多项为 $0$。
\end{minipage}\par\vfill
\noindent\begin{minipage}{\textwidth}\ProblemHeading{例题}{1.22}
设 $a_n=\underbrace{\sqrt{2+\sqrt{2+\cdots+\sqrt2}}}_{n\text{重根式}}$，求 $\displaystyle\lim_{n\to\infty}a_n$。
\end{minipage}\par\vfill
\noindent\begin{minipage}{\textwidth}\ProblemHeading{例题}{1.23}
求 $\displaystyle\lim_{n\to\infty}\left(1-\frac1n\right)^{-n}$。
\end{minipage}\par\vfill
\noindent\begin{minipage}{\textwidth}\ProblemHeading{例题}{1.24}
设 $a_n=\dfrac{\sin1}{1^2}+\dfrac{\sin2}{2^2}+\cdots+\dfrac{\sin n}{n^2}$，证明 $\{a_n\}$ 收敛。
\end{minipage}\par\vfill
\newpage
\noindent\begin{minipage}{\textwidth}\ProblemHeading{例题}{1.25}
设 $a_n=1+\dfrac12+\dfrac13+\cdots+\dfrac1n$，求证 $\{a_n\}$ 发散。
\end{minipage}\par\vfill
\noindent\begin{minipage}{\textwidth}\ProblemHeading{例题}{1.26}
求证若 $\displaystyle\lim_{n\to\infty}a_n=a$，则\[\lim_{n\to\infty}\frac{a_1+a_2+\cdots+a_n}{n}=a.\]
\end{minipage}\par\vfill
\noindent\begin{minipage}{\textwidth}\ProblemHeading{例题}{1.27}
设 $\displaystyle\lim_{n\to\infty}a_n=a$，$\displaystyle\lim_{n\to\infty}b_n=b$，证明\[\lim_{n\to\infty}\frac{a_1b_n+a_2b_{n-1}+\cdots+a_nb_1}{n}=ab.\]
\end{minipage}\par\vfill
\noindent\begin{minipage}{\textwidth}\ProblemHeading{习题}{1.28}
用定义证明下面的结论：
\SubQ{1}{$\displaystyle\lim_{n\to\infty}\frac{n}{5+3n}=\frac13$。}
\SubQ{2}{$\displaystyle\lim_{n\to\infty}\frac{\sin n}{n}=0$。}
\SubQ{3}{$\displaystyle\lim_{n\to\infty}\frac{(-1)^n}{\sqrt{n+1}}=0$。}
\end{minipage}\par\vfill
\newpage
\noindent\begin{minipage}{\textwidth}\ProblemHeading{习题}{1.29}
用定义证明下面的结论：
\SubQ{1}{$\displaystyle\lim_{n\to\infty}\frac{n!}{n^n}=0$。}
\end{minipage}\par\vfill
\noindent\begin{minipage}{\textwidth}\ProblemHeading{习题}{1.30}
证明若数列 $\{a_n\}\ (n\ge1)$ 满足条件：任给正数 $\varepsilon$，存在自然数 $N$，使得当 $n>N$ 时，有 $|a_n-a|<M\varepsilon$（其中 $M$ 为常数），则 $\{a_n\}$ 必以 $a$ 为极限。
\end{minipage}\par\vfill
\noindent\begin{minipage}{\textwidth}\ProblemHeading{习题}{1.31}
证明：当且仅当 $\displaystyle\lim_{n\to\infty}(a_n-a)=0$ 时，有 $\displaystyle\lim_{n\to\infty}a_n=a$。
\end{minipage}\par\vfill
\noindent\begin{minipage}{\textwidth}\ProblemHeading{习题}{1.32}
证明：若 $\displaystyle\lim_{n\to\infty}a_n=a$，则 $\displaystyle\lim_{n\to\infty}|a_n|=|a|$；反之不一定成立（试举例说明）。但若 $\displaystyle\lim_{n\to\infty}|a_n|=0$，则有 $\displaystyle\lim_{n\to\infty}a_n=0$。
\end{minipage}\par\vfill
\newpage
\noindent\begin{minipage}{\textwidth}\ProblemHeading{习题}{1.33}
证明：若 $\displaystyle\lim_{n\to\infty}a_n=0$，又 $|b_n|\le M\ (n=1,2,\ldots)$，则 $\displaystyle\lim_{n\to\infty}a_nb_n=0$。
\end{minipage}\par\vfill
\noindent\begin{minipage}{\textwidth}\ProblemHeading{习题}{1.34}
证明：若数列 $\{a_n\}$ 满足 $\displaystyle\lim_{k\to\infty}a_{2k+1}=a$ 及 $\displaystyle\lim_{k\to\infty}a_{2k}=a$，则 $\displaystyle\lim_{n\to\infty}a_n=a$。
\end{minipage}\par\vfill
\noindent\begin{minipage}{\textwidth}\ProblemHeading{习题}{1.35}
证明下列数列不收敛：
\SubQ{1}{$a_n=(-1)^n\dfrac{n}{n+1}$。}
\SubQ{2}{$a_n=5\left(1-\dfrac2n\right)+(-1)^n$。}
\end{minipage}\par\vfill
\noindent\begin{minipage}{\textwidth}\ProblemHeading{习题}{1.36}
求下列数列的极限：
\SubQ{1}{$a_n=\dfrac{4n^2+5n+2}{3n^2+2n+1}$。}
\SubQ{2}{$a_n=\dfrac1{1\cdot2}+\dfrac1{2\cdot3}+\cdots+\dfrac1{(n-1)n}$。}
\SubQ{3}{$a_n=\left(1-\dfrac13\right)\left(1-\dfrac16\right)\cdots\left(1-\dfrac1{n(n+1)/2}\right),\quad n=2,3,\ldots$。}
\end{minipage}\par\vfill
\newpage
\noindent\begin{minipage}{\textwidth}\ProblemHeading{习题}{1.37}
求下列数列的极限：
\SubQ{1}{$a_n=\left(1-\dfrac1{2^2}\right)\left(1-\dfrac1{3^2}\right)\cdots\left(1-\dfrac1{n^2}\right)$。}
\SubQ{2}{$a_n=(1+q)(1+q^2)(1+q^4)\cdots(1+q^{2^n}),\quad |q|<1$。}
\end{minipage}\par\vfill
\noindent\begin{minipage}{\textwidth}\ProblemHeading{习题}{1.38}
若 $a_n\ne0\ (n=1,2,\ldots)$ 且 $\displaystyle\lim_{n\to\infty}a_n=a$，能否断定 $\displaystyle\lim_{n\to\infty}\frac{a_n}{a_{n+1}}=1$？
\end{minipage}\par\vfill
\noindent\begin{minipage}{\textwidth}\ProblemHeading{习题}{1.39}
若数列 $\{a_n\}$、$\{b_n\}$ 满足 $\displaystyle\lim_{n\to\infty}(a_nb_n)=0$，是否必有 $\displaystyle\lim_{n\to\infty}a_n=0$ 或 $\displaystyle\lim_{n\to\infty}b_n=0$？若还假设 $\displaystyle\lim_{n\to\infty}a_n=a$，回答同样的问题。
\end{minipage}\par\vfill
\noindent\begin{minipage}{\textwidth}\ProblemHeading{习题}{1.40}
若数列 $\{a_n\}$ 收敛，数列 $\{b_n\}$ 发散，则数列 $\{a_n\pm b_n\}$、$\{a_nb_n\}$ 的敛散性如何？举例说明。若数列 $\{a_n\}$ 与 $\{b_n\}$ 皆发散，回答同样的问题。
\end{minipage}\par\vfill
\newpage
\noindent\begin{minipage}{\textwidth}\ProblemHeading{习题}{1.41}
下面的推理是否正确？
\SubQ{1}{设数列 $\{a_n\}$：$a_1=1$，$a_{n+1}=2a_n-1\ (n=1,2,3,\ldots)$，求 $\displaystyle\lim_{n\to\infty}a_n$。所给推理：设 $\displaystyle\lim_{n\to\infty}a_n=a$，在 $a_{n+1}=2a_n-1$ 两边取极限，得 $a=2a-1$，即 $a=1$。}
\SubQ{2}{所给推理为\[\begin{aligned}&\lim_{n\to\infty}\left(\frac1{\sqrt{n^2+1}}+\frac1{\sqrt{n^2+2}}+\cdots+\frac1{\sqrt{n^2+n}}\right)\\&=\lim_{n\to\infty}\frac1{\sqrt{n^2+1}}+\lim_{n\to\infty}\frac1{\sqrt{n^2+2}}+\cdots+\lim_{n\to\infty}\frac1{\sqrt{n^2+n}}\\&=\underbrace{0+0+\cdots+0}_{n\text{个}}=0.\end{aligned}\]}
\SubQ{3}{所给推理为\[\lim_{n\to\infty}\left(1+\frac1n\right)^n=\left(\lim_{n\to\infty}\left(1+\frac1n\right)\right)^n=1^n=1.\]}
\end{minipage}\par\vfill
\newpage
\noindent\begin{minipage}{\textwidth}\ProblemHeading{习题}{1.42}
设数列 $\{a_n\}$ 与 $\{b_n\}$ 分别收敛于 $a,b$。证明：若 $a>b$，则从某一项开始，有 $a_n>b_n$；反过来，若从某项开始恒有 $a_n\ge b_n$，则 $a\ge b$。
\end{minipage}\par\vfill
\noindent\begin{minipage}{\textwidth}\ProblemHeading{习题}{1.43}
设数列 $\{a_n\}$、$\{b_n\}$ 分别收敛于 $a$ 及 $b$。记 $c_n=\max\{a_n,b_n\}$、$d_n=\min\{a_n,b_n\}\ (n=1,2,\ldots)$。证明：\[\lim_{n\to\infty}c_n=\max\{a,b\},\qquad\lim_{n\to\infty}d_n=\min\{a,b\}.\]
\end{minipage}\par\vfill
\noindent\begin{minipage}{\textwidth}\ProblemHeading{习题}{1.44}
求下列极限：
\SubQ{1}{$\displaystyle\lim_{n\to\infty}\left(\frac1{(n+1)^2}+\frac1{(n+2)^2}+\cdots+\frac1{(2n)^2}\right)$。}
\SubQ{2}{$\displaystyle\lim_{n\to\infty}[(n+1)^k-n^k]$，其中 $0<k<1$。}
\SubQ{3}{$\displaystyle\lim_{n\to\infty}\left(\sqrt2\cdot\sqrt[4]2\cdot\sqrt[8]2\cdots\sqrt[2^n]2\right)$。}
\end{minipage}\par\vfill
\noindent\begin{minipage}{\textwidth}\ProblemHeading{习题}{1.45}
求下列极限：
\SubQ{1}{$\displaystyle\lim_{n\to\infty}\sqrt[n]{n^2-n+2}$。}
\SubQ{2}{$\displaystyle\lim_{n\to\infty}\sqrt[n]{\cos^2 1+\cos^2 2+\cdots+\cos^2 n}$。}
\end{minipage}\par\vfill
\newpage
\noindent\begin{minipage}{\textwidth}\ProblemHeading{习题}{1.46}
设 $a_1,\ldots,a_m$ 为 $m$ 个正数，证明：\[\lim_{n\to\infty}\sqrt[n]{a_1^n+a_2^n+\cdots+a_m^n}=\max\{a_1,a_2,\ldots,a_m\}.\]
\end{minipage}\par\vfill
\noindent\begin{minipage}{\textwidth}\ProblemHeading{习题}{1.47}
证明下列数列收敛：
\SubQ{1}{$a_n=\left(1-\dfrac12\right)\left(1-\dfrac1{2^2}\right)\cdots\left(1-\dfrac1{2^n}\right)$。}
\SubQ{2}{$a_n=\dfrac1{3+1}+\dfrac1{3^2+1}+\cdots+\dfrac1{3^n+1}$。}
\SubQ{3}{$a_n=\alpha_0+\alpha_1q+\cdots+\alpha_nq^n$，其中 $|\alpha_k|\le M\ (k=1,2,\ldots)$，而 $|q|<1$。}
\end{minipage}\par\vfill
\noindent\begin{minipage}{\textwidth}\ProblemHeading{习题}{1.48}
证明下列数列收敛：
\SubQ{1}{$a_n=\dfrac{\cos1}{1\cdot2}+\dfrac{\cos2}{2\cdot3}+\dfrac{\cos3}{3\cdot4}+\cdots+\dfrac{\cos n}{n(n+1)}$。}
\end{minipage}\par\vfill
\noindent\begin{minipage}{\textwidth}\ProblemHeading{习题}{1.49}
证明下列数列收敛，并求出其极限：
\SubQ{1}{$a_n=\dfrac{n}{c^n}\ (c>1)$。}
\SubQ{2}{$a_1=\dfrac c2$，$a_{n+1}=\dfrac c2+\dfrac{a_n^2}{2}\ (0\le c\le1)$。}
\SubQ{3}{$a>0$，$a_0>0$，$a_{n+1}=\dfrac12\left(a_n+\dfrac a{a_n}\right)$。}
\end{minipage}\par\vfill
\newpage
\noindent\begin{minipage}{\textwidth}\ProblemHeading{习题}{1.50}
证明下列数列收敛，并求出其极限：
\SubQ{1}{$a_0=1$，$a_n=1+\dfrac{a_{n-1}}{a_{n-1}+1}$。}
\SubQ{2}{$a_n=\underbrace{\sin\sin\cdots\sin}_{n\text{个}\sin}1$。}
\end{minipage}\par\vfill
\noindent\begin{minipage}{\textwidth}\ProblemHeading{习题}{1.51}
设 $a_n\le a\le b_n\ (n=1,2,\ldots)$，且 $\displaystyle\lim_{n\to\infty}(a_n-b_n)=0$，求证：$\displaystyle\lim_{n\to\infty}a_n=a$，$\displaystyle\lim_{n\to\infty}b_n=a$。
\end{minipage}\par\vfill
\noindent\begin{minipage}{\textwidth}\ProblemHeading{习题}{1.52}
证明：若 $a_n>0$，且 $\displaystyle\lim_{n\to\infty}\frac{a_n}{a_{n+1}}=l>1$，则 $\displaystyle\lim_{n\to\infty}a_n=0$。
\end{minipage}\par\vfill
\noindent\begin{minipage}{\textwidth}\ProblemHeading{习题}{1.53}
设 $\{a_n\}$、$\{b_n\}$ 是正数列，满足 $\dfrac{a_{n+1}}{a_n}\le\dfrac{b_{n+1}}{b_n}$，$n=1,2,\ldots$。求证：若 $\{b_n\}$ 收敛，则 $\{a_n\}$ 收敛。
\end{minipage}\par\vfill
\newpage
\noindent\begin{minipage}{\textwidth}\ProblemHeading{习题}{1.54}
利用极限 $\displaystyle\lim_{n\to\infty}\left(1+\frac1n\right)^n=e$，求下列数列的极限：
\SubQ{1}{$a_n=\left(1+\dfrac1{2n+1}\right)^{2n+1}$。}
\SubQ{2}{$a_n=\left(1-\dfrac1{n-2}\right)^{n+1}$。}
\SubQ{3}{$a_n=\left(\dfrac{1+n}{2+n}\right)^n$。}
\end{minipage}\par\vfill
\noindent\begin{minipage}{\textwidth}\ProblemHeading{习题}{1.55}
利用极限 $\displaystyle\lim_{n\to\infty}\left(1+\frac1n\right)^n=e$，求下列数列的极限：
\SubQ{1}{$a_n=\left(1+\dfrac1{n^3}\right)^{2n^3}$。}
\end{minipage}\par\vfill
\noindent\begin{minipage}{\textwidth}\ProblemHeading{习题}{1.56}
设 $\displaystyle\lim_{n\to\infty}a_n=\infty$，且 $|b_n|\ge b>0\ (n=1,2,\ldots)$，证明：$\displaystyle\lim_{n\to\infty}a_nb_n=\infty$。
\end{minipage}\par\vfill
\noindent\begin{minipage}{\textwidth}\ProblemHeading{习题}{1.57}
确定 $n\to\infty$ 时，$\sqrt[n]{n!}$ 与 $n\sin\dfrac{n\pi}{2}\ (n\ge1)$ 是否有界，是否趋于无穷大。
\end{minipage}\par\vfill
\newpage
\noindent\begin{minipage}{\textwidth}\ProblemHeading{习题}{1.58}
设数列 $\{a_n\}$ 由 $a_1=1$，$a_{n+1}=a_n+\dfrac1{a_n}\ (n\ge1)$ 定义，证明：$a_n\to+\infty\ (n\to\infty)$。
\end{minipage}\par\vfill
\noindent\begin{minipage}{\textwidth}\ProblemHeading{习题}{1.59}
给出 $0/0$ 型 Stolz 定理的证明：设 $\{a_n\}$、$\{b_n\}$ 是两个收敛于 $0$ 的数列，且 $\{b_n\}$ 是严格递减数列。如果\[\lim_{n\to\infty}\frac{a_{n+1}-a_n}{b_{n+1}-b_n}=A,\]则 $\displaystyle\lim_{n\to\infty}\frac{a_n}{b_n}=A$，其中 $A$ 可以是实数，也可以是 $+\infty$ 或 $-\infty$。
\end{minipage}\par\vfill
\noindent\begin{minipage}{\textwidth}\ProblemHeading{例题}{1.60（改写）}
设 $n$ 为非负整数，$a_0,\ldots,a_n$ 为实数，$f(x)=a_nx^n+\cdots+a_1x+a_0$，$g(x)$ 为非零多项式。求多项式函数 $f(x)$ 及有理函数 $f(x)/g(x)$ 的定义域。
\end{minipage}\par\vfill
\noindent\begin{minipage}{\textwidth}\ProblemHeading{例题}{1.61（改写）}
考虑幂函数 $f(x)=x^\alpha$，完成下列问题。
\SubQ{1}{当 $\alpha$ 为整数时，证明 $f(x)$ 是有理函数。}
\SubQ{2}{当 $\alpha=1/m$，其中 $m\ge2$ 为整数时，分别讨论 $m$ 为奇数、偶数时 $f$ 的定义域。}
\SubQ{3}{当 $\alpha$ 为无理数时，按通常实数幂的约定，写出 $f$ 的定义域。}
\end{minipage}\par\vfill
\newpage
\noindent\begin{minipage}{\textwidth}\ProblemHeading{例题}{1.62（改写）}
双曲函数定义为 $\sinh x=(e^x-e^{-x})/2$、$\cosh x=(e^x+e^{-x})/2$、$\tanh x=(e^x-e^{-x})/(e^x+e^{-x})$ 及 $\coth x=(e^x+e^{-x})/(e^x-e^{-x})$（最后一式要求 $x\ne0$）。证明下列恒等式。
\SubQ{1}{$\sinh(-x)=-\sinh x$。}
\SubQ{2}{$\cosh(-x)=\cosh x$。}
\SubQ{3}{$\sinh(x\pm y)=\sinh x\cosh y\pm\cosh x\sinh y$。}
\end{minipage}\par\vfill
\noindent\begin{minipage}{\textwidth}\ProblemHeading{例题}{1.63（改写）}
双曲函数定义为 $\sinh x=(e^x-e^{-x})/2$、$\cosh x=(e^x+e^{-x})/2$、$\tanh x=(e^x-e^{-x})/(e^x+e^{-x})$ 及 $\coth x=(e^x+e^{-x})/(e^x-e^{-x})$（最后一式要求 $x\ne0$）。证明下列恒等式。
\SubQ{1}{$\cosh(x\pm y)=\cosh x\cosh y\pm\sinh x\sinh y$。}
\SubQ{2}{$\cosh^2x-\sinh^2x=1$。}
\SubQ{3}{$\sinh2x=2\sinh x\cosh x$。}
\end{minipage}\par\vfill
\noindent\begin{minipage}{\textwidth}\ProblemHeading{例题}{1.64（改写）}
双曲函数定义为 $\sinh x=(e^x-e^{-x})/2$、$\cosh x=(e^x+e^{-x})/2$、$\tanh x=(e^x-e^{-x})/(e^x+e^{-x})$ 及 $\coth x=(e^x+e^{-x})/(e^x-e^{-x})$（最后一式要求 $x\ne0$）。证明下列恒等式。
\SubQ{1}{$\cosh2x=\cosh^2x+\sinh^2x$。}
\end{minipage}\par\vfill
\noindent\begin{minipage}{\textwidth}\ProblemHeading{例题}{1.65}
求双曲正弦函数 $y=\sinh x=(e^x-e^{-x})/2$ 的反函数。
\end{minipage}\par\vfill
\newpage
\noindent\begin{minipage}{\textwidth}\ProblemHeading{例题}{1.66}
设 $k$ 是正整数，证明 $\displaystyle\lim_{x\to\infty}\frac1{x^k}=0$。
\end{minipage}\par\vfill
\noindent\begin{minipage}{\textwidth}\ProblemHeading{例题}{1.67}
设 $0<a<1$，证明 $\displaystyle\lim_{x\to+\infty}a^x=0$。
\end{minipage}\par\vfill
\noindent\begin{minipage}{\textwidth}\ProblemHeading{例题}{1.68}
证明 $\displaystyle\lim_{x\to-\infty}e^x=0$。
\end{minipage}\par\vfill
\noindent\begin{minipage}{\textwidth}\ProblemHeading{例题}{1.69}
证明下列极限。
\SubQ{1}{$\displaystyle\lim_{x\to-\infty}\arctan x=-\frac\pi2$。}
\SubQ{2}{$\displaystyle\lim_{x\to+\infty}\arctan x=\frac\pi2$。}
\end{minipage}\par\vfill
\newpage
\noindent\begin{minipage}{\textwidth}\ProblemHeading{例题}{1.70}
设 $f(x)=x$，证明 $\displaystyle\lim_{x\to x_0}f(x)=x_0$。
\end{minipage}\par\vfill
\noindent\begin{minipage}{\textwidth}\ProblemHeading{例题}{1.71}
求 $\displaystyle\lim_{x\to0}x\sin\frac1x$。
\end{minipage}\par\vfill
\noindent\begin{minipage}{\textwidth}\ProblemHeading{例题}{1.72}
证明 $\displaystyle\lim_{x\to1}\frac{x^2-1}{x^2-x}=2$。
\end{minipage}\par\vfill
\noindent\begin{minipage}{\textwidth}\ProblemHeading{例题}{1.73}
设 $a>0$，证明 $\displaystyle\lim_{x\to0}a^x=1$。
\end{minipage}\par\vfill
\newpage
\noindent\begin{minipage}{\textwidth}\ProblemHeading{例题}{1.74}
设 $\displaystyle f(x)=\begin{cases}3x+1,&x\ge1,\\2x-5,&x<1.\end{cases}$ 证明 $\displaystyle\lim_{x\to1}f(x)$ 不存在。
\end{minipage}\par\vfill
\noindent\begin{minipage}{\textwidth}\ProblemHeading{例题}{1.75}
求 $\displaystyle\lim_{x\to\infty}\frac{x^2+2x+3}{x^2+3x+1}$。
\end{minipage}\par\vfill
\noindent\begin{minipage}{\textwidth}\ProblemHeading{例题}{1.76}
设 $P(x)=a_nx^n+a_{n-1}x^{n-1}+\cdots+a_1x+a_0$，证明对于任意一点 $x_0$，有 $\displaystyle\lim_{x\to x_0}P(x)=P(x_0)$。
\end{minipage}\par\vfill
\noindent\begin{minipage}{\textwidth}\ProblemHeading{例题}{1.77}
求 $\displaystyle\lim_{x\to-1}\left(\frac1{x+1}-\frac3{x^3+1}\right)$。
\end{minipage}\par\vfill
\newpage
\noindent\begin{minipage}{\textwidth}\ProblemHeading{例题}{1.78}
设 $a>0$，证明 $\displaystyle\lim_{x\to x_0}a^x=a^{x_0}$。
\end{minipage}\par\vfill
\noindent\begin{minipage}{\textwidth}\ProblemHeading{例题}{1.79}
证明当 $x\to0$ 时，$f(x)=\sin(1/x)$ 没有极限。
\end{minipage}\par\vfill
\noindent\begin{minipage}{\textwidth}\ProblemHeading{例题}{1.80}
证明下列极限。
\SubQ{1}{$\displaystyle\lim_{x\to0}\cos x=1$。}
\SubQ{2}{$\displaystyle\lim_{x\to0}\frac{\tan x}{x}=1$。}
\end{minipage}\par\vfill
\noindent\begin{minipage}{\textwidth}\ProblemHeading{例题}{1.81}
求 $\displaystyle\lim_{x\to0}\frac{1-\cos x}{x^2}$。
\end{minipage}\par\vfill
\newpage
\noindent\begin{minipage}{\textwidth}\ProblemHeading{例题}{1.82}
设 $x_0$ 是任意实数，证明 $\displaystyle\lim_{x\to x_0}\sin x=\sin x_0$。
\end{minipage}\par\vfill
\noindent\begin{minipage}{\textwidth}\ProblemHeading{例题}{1.83（改写）}
证明 $x^2+3x+5\sim x^2$（$x\to\infty$）。
\end{minipage}\par\vfill
\noindent\begin{minipage}{\textwidth}\ProblemHeading{例题}{1.84}
证明 $\ln x=o(x)$（$x\to+\infty$）。
\end{minipage}\par\vfill
\noindent\begin{minipage}{\textwidth}\ProblemHeading{例题}{1.85}
求 $\displaystyle\lim_{x\to0}\frac{\sin x}{\sqrt{1+x}-1}$。
\end{minipage}\par\vfill
\newpage
\noindent\begin{minipage}{\textwidth}\ProblemHeading{例题}{1.86}
求 $\displaystyle\lim_{x\to0}\frac{\tan x-\sin x}{x^3}$。
\end{minipage}\par\vfill
\noindent\begin{minipage}{\textwidth}\ProblemHeading{习题}{1.87}
按定义证明下列极限。
\SubQ{1}{$\displaystyle\lim_{x\to-\infty}a^x=0$（$a>1$）。}
\SubQ{2}{$\displaystyle\lim_{x\to\infty}\frac{x-1}{x+1}=1$。}
\SubQ{3}{$\displaystyle\lim_{x\to-1}\frac{x^2-1}{x^2+x}=2$。}
\end{minipage}\par\vfill
\noindent\begin{minipage}{\textwidth}\ProblemHeading{习题}{1.88}
按定义证明下列极限。
\SubQ{1}{$\displaystyle\lim_{x\to0^+}x^{1/q}=0$（$q$ 为正整数）。}
\end{minipage}\par\vfill
\noindent\begin{minipage}{\textwidth}\ProblemHeading{习题}{1.89}
求下列极限。
\SubQ{1}{$\displaystyle\lim_{x\to1}\left(x^5-5x+2+\frac1x\right)$。}
\SubQ{2}{$\displaystyle\lim_{x\to1}\frac{x^n-1}{x-1}$（$n$ 为自然数）。}
\SubQ{3}{$\displaystyle\lim_{x\to1}\frac{x^2-1}{2x^2-x-1}$。}
\end{minipage}\par\vfill
\newpage
\noindent\begin{minipage}{\textwidth}\ProblemHeading{习题}{1.90}
求下列极限。
\SubQ{1}{$\displaystyle\lim_{x\to\infty}\frac{(3x+6)^{70}(8x-5)^{20}}{(5x-1)^{90}}$。}
\end{minipage}\par\vfill
\noindent\begin{minipage}{\textwidth}\ProblemHeading{习题}{1.91}
证明下列极限不存在。
\SubQ{1}{$\displaystyle\lim_{x\to+\infty}\sin x$。}
\SubQ{2}{$\displaystyle\lim_{x\to0}\frac{|x|}{x}$。}
\end{minipage}\par\vfill
\noindent\begin{minipage}{\textwidth}\ProblemHeading{习题}{1.92}
设函数 $f(x)$ 在正无穷大处的极限为 $l$，证明：对于任意趋于正无穷大的数列 $\{a_n\}$，有 $\displaystyle\lim_{n\to\infty}f(a_n)=l$。特别地，$\displaystyle\lim_{n\to\infty}f(n)=l$。
\end{minipage}\par\vfill
\noindent\begin{minipage}{\textwidth}\ProblemHeading{习题}{1.93}
讨论下列函数在 $x=0$ 处的极限（$\lfloor x\rfloor$ 表示不超过 $x$ 的最大整数）。
\SubQ{1}{$f(x)=\lfloor x\rfloor$。}
\SubQ{2}{$f(x)=\operatorname{sgn}x$。}
\SubQ{3}{$\displaystyle f(x)=\begin{cases}2^x,&x>0,\\0,&x=0,\\1+x^2,&x<0.\end{cases}$}
\end{minipage}\par\vfill
\newpage
\noindent\begin{minipage}{\textwidth}\ProblemHeading{习题}{1.94}
讨论下列函数在 $x=0$ 处的极限（$\lfloor x\rfloor$ 表示不超过 $x$ 的最大整数）。
\SubQ{1}{$\displaystyle f(x)=\begin{cases}\cos(1/x),&x>0,\\x,&x\le0.\end{cases}$}
\end{minipage}\par\vfill
\noindent\begin{minipage}{\textwidth}\ProblemHeading{习题}{1.95}
求 $\displaystyle\lim_{n\to\infty}\cos\frac x2\cos\frac x{2^2}\cdots\cos\frac x{2^n}$。
\end{minipage}\par\vfill
\noindent\begin{minipage}{\textwidth}\ProblemHeading{习题}{1.96}
求证：$\displaystyle\lim_{n\to\infty}\left(\sin\frac\alpha{n^2}+\sin\frac{2\alpha}{n^2}+\cdots+\sin\frac{n\alpha}{n^2}\right)=\frac\alpha2$。
\end{minipage}\par\vfill
\noindent\begin{minipage}{\textwidth}\ProblemHeading{习题}{1.97}
证明：$\displaystyle\lim_{x\to\infty}f(x)=l$ 与 $\displaystyle\lim_{x\to0}f(1/x)=l$ 等价。叙述并证明 $x\to+\infty$ 及 $x\to-\infty$ 时类似的结论。
\end{minipage}\par\vfill
\newpage
\noindent\begin{minipage}{\textwidth}\ProblemHeading{习题}{1.98}
求下列极限。
\SubQ{1}{$\displaystyle\lim_{x\to0}\frac{\tan2x}{\sin5x}$。}
\SubQ{2}{$\displaystyle\lim_{x\to0}\frac{\cos x-\cos3x}{x^2}$。}
\SubQ{3}{$\displaystyle\lim_{x\to+\infty}\left(\frac{x+1}{2x-1}\right)^x$。}
\end{minipage}\par\vfill
\noindent\begin{minipage}{\textwidth}\ProblemHeading{习题}{1.99}
求下列极限。
\SubQ{1}{$\displaystyle\lim_{x\to\infty}\left(\frac{x^2+1}{x^2-1}\right)^{x^2}$。}
\end{minipage}\par\vfill
\noindent\begin{minipage}{\textwidth}\ProblemHeading{习题}{1.100}
求下列极限。
\SubQ{1}{$\displaystyle\lim_{x\to+\infty}\frac{\arctan x}{x}$。}
\SubQ{2}{$\displaystyle\lim_{x\to0}x^2\sin\frac1x$。}
\SubQ{3}{$\displaystyle\lim_{x\to2}\frac{x^3-2x^2}{x-2}$。}
\end{minipage}\par\vfill
\noindent\begin{minipage}{\textwidth}\ProblemHeading{习题}{1.101}
求下列极限。
\SubQ{1}{$\displaystyle\lim_{x\to\infty}(2x^2-x+1)$。}
\end{minipage}\par\vfill
\newpage
\noindent\begin{minipage}{\textwidth}\ProblemHeading{习题}{1.102}
按定义证明下列极限。
\SubQ{1}{$\displaystyle\lim_{x\to+\infty}\log_a x=+\infty$（$a>1$）。}
\SubQ{2}{$\displaystyle\lim_{x\to0^+}\log_a x=-\infty$（$a>1$）。}
\SubQ{3}{$\displaystyle\lim_{x\to(\pi/2)^-}\tan x=+\infty$。}
\end{minipage}\par\vfill
\noindent\begin{minipage}{\textwidth}\ProblemHeading{习题}{1.103}
按定义证明下列极限。
\SubQ{1}{$\displaystyle\lim_{x\to0^+}e^{1/x}=+\infty$。}
\end{minipage}\par\vfill
\noindent\begin{minipage}{\textwidth}\ProblemHeading{习题}{1.104}
证明：函数 $y=x\sin x$ 在 $(0,+\infty)$ 内无界，但当 $x\to+\infty$ 时，这个函数并不是无穷大量。
\end{minipage}\par\vfill
\noindent\begin{minipage}{\textwidth}\ProblemHeading{习题}{1.105}
函数 $y=\dfrac1x\cos\dfrac1x$ 在区间 $(0,1)$ 内是否有界？当 $x\to0^+$ 时，这个函数是否为无穷大量？
\end{minipage}\par\vfill
\newpage
\noindent\begin{minipage}{\textwidth}\ProblemHeading{习题}{1.106}
记函数 $y=f(x)$ 的曲线为 $C$。若动点沿 $C$ 无限远离原点时到某固定直线的距离趋于零，则称该直线为 $C$ 的渐近线。垂直渐近线 $x=x_0$ 的判别条件为 $f(x_0-0)=\infty$ 或 $f(x_0+0)=\infty$；水平渐近线 $y=b$ 的判别条件为 $\lim_{x\to+\infty}f(x)=b$ 或 $\lim_{x\to-\infty}f(x)=b$。完成下列问题。
\SubQ{1}{证明：直线 $y=ax+b$（$a\ne0$）为 $C$ 的斜渐近线的充分必要条件是 $a=\lim_{x\to+\infty}f(x)/x$、$b=\lim_{x\to+\infty}(f(x)-ax)$，或这两个式子中的极限过程同时改为 $x\to-\infty$（所说的极限均存在）。}
\SubQ{2}{求曲线 $y=x\ln\left(e+\dfrac1x\right)$ 的渐近线方程。}
\SubQ{3}{求曲线 $y=\dfrac{3x^2-2x+3}{x-1}$ 的渐近线方程。}
\end{minipage}\par\vfill
\newpage
\noindent\begin{minipage}{\textwidth}\ProblemHeading{习题}{1.107}
证明同一极限过程中等价无穷小量具有下列性质；所涉及的无穷小量在该过程中均不取零值。
\SubQ{1}{自反性：$\alpha(x)\sim\alpha(x)$。}
\SubQ{2}{对称性：若 $\alpha(x)\sim\beta(x)$，则 $\beta(x)\sim\alpha(x)$。}
\SubQ{3}{传递性：若 $\alpha(x)\sim\beta(x)$，$\beta(x)\sim\gamma(x)$，则 $\alpha(x)\sim\gamma(x)$。}
\end{minipage}\par\vfill
\noindent\begin{minipage}{\textwidth}\ProblemHeading{习题}{1.108}
当 $x\to0$ 时，比较下列无穷小量的阶。
\SubQ{1}{$\tan x-\sin x$ 与 $x^3$。}
\SubQ{2}{$x^3+x^2$ 与 $\sin^2x$。}
\SubQ{3}{$1-\cos x$ 与 $x^2$。}
\end{minipage}\par\vfill
\noindent\begin{minipage}{\textwidth}\ProblemHeading{习题}{1.109}
当 $x\to+\infty$ 时，比较下列无穷大量的阶。
\SubQ{1}{$n$ 次多项式 $P_n(x)$ 与 $m$ 次多项式 $P_m(x)$（$m,n$ 均为正整数）。}
\SubQ{2}{$x^\alpha$ 与 $x^\beta$（$\alpha,\beta>0$）。}
\SubQ{3}{$a^x$ 与 $b^x$（$a,b>1$）。}
\end{minipage}\par\vfill
\noindent\begin{minipage}{\textwidth}\ProblemHeading{习题}{1.110}
试用等价无穷小量代换的方法计算下列极限。
\SubQ{1}{$\displaystyle\lim_{x\to0}\frac{\sin mx}{\sin nx}$（$m,n$ 均为正整数）。}
\SubQ{2}{$\displaystyle\lim_{x\to0}\frac{\tan ax}{x}$。}
\SubQ{3}{$\displaystyle\lim_{x\to0}\frac{\sqrt[n]{1+\sin x}-1}{\arctan x}$。}
\end{minipage}\par\vfill
\newpage
\noindent\begin{minipage}{\textwidth}\ProblemHeading{习题}{1.111}
试用等价无穷小量代换的方法计算下列极限。
\SubQ{1}{$\displaystyle\lim_{x\to0}\frac{\sqrt2-\sqrt{1+\cos x}}{\sin^2x}$。}
\SubQ{2}{$\displaystyle\lim_{x\to0}\frac{\sqrt{1+x+x^2}-1}{\sin2x}$。}
\SubQ{3}{$\displaystyle\lim_{x\to0}\frac{\sqrt{1+x^2}-1}{1-\cos x}$。}
\end{minipage}\par\vfill
\noindent\begin{minipage}{\textwidth}\ProblemHeading{习题}{1.112}
求下列数列的极限。
\SubQ{1}{$a_n=\dfrac12\cdot\dfrac34\cdots\dfrac{2n-1}{2n}$。}
\SubQ{2}{$a_n=\dfrac{10}1\cdot\dfrac{11}3\cdot\dfrac{12}5\cdots\dfrac{n+9}{2n-1}$。}
\SubQ{3}{$a_1>1$，$a_{n+1}=2-1/a_n$，$n=1,2,\ldots$。}
\end{minipage}\par\vfill
\noindent\begin{minipage}{\textwidth}\ProblemHeading{习题}{1.113}
求下列数列的极限。
\SubQ{1}{$a_1=3$，$a_{n+1}=1/(1+a_n)$，$n=1,2,\ldots$。}
\end{minipage}\par\vfill
\noindent\begin{minipage}{\textwidth}\ProblemHeading{习题}{1.114}
设 $\{a_n\}$ 为单调递增数列，并且收敛于 $a$，证明：对一切 $n$ 有 $a_n\le a$。对单调递减且有极限的数列，叙述并证明类似的结论。
\end{minipage}\par\vfill
\newpage
\noindent\begin{minipage}{\textwidth}\ProblemHeading{习题}{1.115}
证明下面的数列收敛。
\SubQ{1}{$a_n=1+\dfrac1{2^2}+\cdots+\dfrac1{n^2}$。}
\SubQ{2}{$a_n=\left(1+\dfrac12\right)\left(1+\dfrac1{2^2}\right)\cdots\left(1+\dfrac1{2^n}\right)$。}
\end{minipage}\par\vfill
\noindent\begin{minipage}{\textwidth}\ProblemHeading{习题}{1.116}
试构造一个发散数列 $\{a_n\}$，满足条件：对任意正数 $\varepsilon$，存在自然数 $N$，使当 $n>N$ 时有 $|a_{n+1}-a_n|<\varepsilon$。
\end{minipage}\par\vfill
\noindent\begin{minipage}{\textwidth}\ProblemHeading{习题}{1.117}
设数列 $\{a_n\}$ 满足：存在常数 $M$，使得对一切 $n$ 有 $A_n=|a_2-a_1|+|a_3-a_2|+\cdots+|a_{n+1}-a_n|\le M$。证明：
\SubQ{1}{数列 $\{A_n\}$ 收敛。}
\SubQ{2}{数列 $\{a_n\}$ 也收敛。}
\end{minipage}\par\vfill
\noindent\begin{minipage}{\textwidth}\ProblemHeading{习题}{1.118}
设 $\{a_n\}$ 是正的严格递增数列。证明：若 $\{a_{n+1}-a_n\}$ 有界，则对任意 $\alpha\in(0,1)$，有 $\lim_{n\to\infty}(a_{n+1}^\alpha-a_n^\alpha)=0$。并说明逆命题不成立，即存在这样的正的严格递增数列，使该极限对任意 $\alpha\in(0,1)$ 均为零，而 $\{a_{n+1}-a_n\}$ 无界。
\end{minipage}\par\vfill
\newpage
\noindent\begin{minipage}{\textwidth}\ProblemHeading{习题}{1.119}
设数列 $\{a_n\}$ 满足 $\lim_{n\to\infty}(a_{n+1}-a_n)=a$，证明 $\displaystyle\lim_{n\to\infty}\frac{a_n}{n}=a$。
\end{minipage}\par\vfill
\noindent\begin{minipage}{\textwidth}\ProblemHeading{习题}{1.120}
证明：若 $\lim_{n\to\infty}a_n=a$，且 $a_n>0$，则 $\displaystyle\lim_{n\to\infty}\sqrt[n]{a_1a_2\cdots a_n}=a$。
\end{minipage}\par\vfill
\noindent\begin{minipage}{\textwidth}\ProblemHeading{习题}{1.121}
证明：若 $a_n>0$，且 $\displaystyle\lim_{n\to\infty}\frac{a_{n+1}}{a_n}$ 存在，则 $\displaystyle\lim_{n\to\infty}\sqrt[n]{a_n}$ 也存在，并且 $\displaystyle\lim_{n\to\infty}\sqrt[n]{a_n}=\lim_{n\to\infty}\frac{a_{n+1}}{a_n}$。
\end{minipage}\par\vfill
\noindent\begin{minipage}{\textwidth}\ProblemHeading{习题}{1.122}
求下列极限。
\SubQ{1}{$\displaystyle\lim_{n\to\infty}\frac{1+\sqrt2+\sqrt[3]3+\cdots+\sqrt[n]n}{n}$。}
\SubQ{2}{$\displaystyle\lim_{n\to\infty}\frac n{\sqrt[n]{n!}}$。}
\end{minipage}\par\vfill
\newpage
\noindent\begin{minipage}{\textwidth}\ProblemHeading{习题}{1.123}
已知 $\lim_{n\to\infty}a_n=a$，求证 $\displaystyle\lim_{n\to\infty}\frac{a_1+2a_2+\cdots+na_n}{n^2}=\frac a2$。
\end{minipage}\par\vfill
\noindent\begin{minipage}{\textwidth}\ProblemHeading{习题}{1.124}
设 $a_n\to a\in\mathbb R$，$\{b_n\}$ 是正数列，$\displaystyle c_n=\frac{a_1b_1+a_2b_2+\cdots+a_nb_n}{b_1+b_2+\cdots+b_n}$。求证：
\SubQ{1}{$\{c_n\}$ 收敛。}
\SubQ{2}{若 $b_1+b_2+\cdots+b_n\to+\infty$，则 $\lim_{n\to\infty}c_n=a$。}
\end{minipage}\par\vfill
\noindent\begin{minipage}{\textwidth}\ProblemHeading{习题}{1.125}
证明：$\displaystyle\lim_{x\to+\infty}\left(1+\frac1{x^p}\right)^x=\begin{cases}1,&p>1,\\e,&p=1,\\+\infty,&p<1.\end{cases}$
\end{minipage}\par\vfill
\noindent\begin{minipage}{\textwidth}\ProblemHeading{习题}{1.126}
设 $f(x)$ 为周期函数，且 $\lim_{x\to\infty}f(x)=0$，证明 $f(x)$ 恒为零。
\end{minipage}\par\vfill
\newpage
\noindent\begin{minipage}{\textwidth}\ProblemHeading{习题}{1.127}
证明下列关于单侧极限的结论。
\SubQ{1}{函数 $f(x)$ 在 $x\to x_0^-$ 时有极限 $l$ 的充分必要条件是：对于任意一个以 $x_0$ 为极限的单调递增数列 $\{a_n\}$（$a_n\ne x_0$），都有 $\lim_{n\to\infty}f(a_n)=l$。}
\SubQ{2}{函数 $f(x)$ 在 $x\to x_0^+$ 时有极限 $l$ 的充分必要条件是：对于任意一个以 $x_0$ 为极限的单调递减数列 $\{a_n\}$（$a_n\ne x_0$），都有 $\lim_{n\to\infty}f(a_n)=l$。}
\end{minipage}\par\vfill
\noindent\begin{minipage}{\textwidth}\ProblemHeading{习题}{1.128}
设 $\xi$ 是无理数，$a,b$ 是实数且 $a<b$。求证：存在整数 $m,n$，使得 $m+n\xi\in(a,b)$，即集合 $S=\{m+n\xi\mid m,n\in\mathbb Z\}$ 在 $\mathbb R$ 中稠密。
\end{minipage}\par\vfill
